\chapter{Prerequisites: The Concepts We Need First}
\label{ch:prereq}

\Cref{ch:bigpicture} gave the map. Before we descend into each box, this chapter
assembles the toolbox. Five ideas recur in every later chapter, namely the
language model that generates, the embeddings that measure meaning, the
similarity search that scales them, the calibration that makes a score
trustworthy, the entailment test that grounds an answer, and the low-rank
adaptation that lets the model learn without forgetting. None of these is unique
to \caem{}, but each is used in a specific way, with specific numbers, that the
rest of the book assumes. We treat each as a self-contained briefing: the general
idea first, then exactly how \caem{} uses it, taken from the live code.

\begin{intuition}[Why a whole chapter on background]
Every component of \caem{} is a familiar tool put to a precise purpose. The
verifier is mostly entailment and calibration; the router is mostly embeddings and
similarity search; the self-improvement loop is mostly low-rank adaptation under a
guard. If the tools are clear, the architecture reads as a series of natural
choices rather than a wall of jargon. So this chapter front-loads the toolbox
once, and the later chapters simply reach for it.
\end{intuition}

\section{The generator: what a decoder language model does}
\label{sec:prereq-llm}

At the centre of \caem{} sits one frozen language model that produces text. It is
worth being precise about what ``produces text'' means, because two of the
verifier's nine signals are read directly off the mechanics below.

A modern generator is an \emph{autoregressive decoder}. It reads a sequence of
tokens, which are sub-word units rather than whole words, and predicts the next
token, one at a time. For each position it produces a vector of \emph{logits},
one real number per token in its vocabulary. A softmax turns those logits into a
probability distribution over the vocabulary, and the model either takes the most
probable token (\emph{greedy decoding}) or samples from the distribution. The
chosen token is appended to the sequence and the process repeats until a stop
token appears.

\begin{precise}[From logits to a per-token probability]
For vocabulary position $v$ at step $t$, the model emits a logit $z_{t,v}$. The
probability of token $v$ is the softmax
\[
  p_t(v) \;=\; \frac{\exp(z_{t,v})}{\sum_{w}\exp(z_{t,w})}.
\]
The probability the model assigned to the token it \emph{actually} produced,
$p_t(\hat v_t)$, is the raw material of the token-level confidence signal. A
sequence the model found easy has high per-token probabilities; a sequence it was
unsure about has low ones. Aggregating these across the generated answer, by a
geometric mean of the per-token probabilities, gives the $\utoken$ signal that
appears in both the pre-routing confidence of \Cref{ch:upre} and the verifier of
\Cref{ch:verifier}.
\end{precise}

Two further ideas matter. \textbf{Temperature} rescales the logits before the
softmax: dividing by a temperature above one flattens the distribution and makes
sampling more random, while dividing by a value below one sharpens it. We will
meet temperature again, in a completely different role, as a \emph{calibration}
knob in \Cref{sec:prereq-calibration}; the two uses share a formula but not a
purpose. \textbf{Chain-of-thought} is the practice of having the model generate
intermediate reasoning before its final answer. \caem{} uses a scaffolded form,
prompting the model to emit a \emph{Reasoning} field and then an \emph{Answer}
field, so that the reasoning is available for inspection and the short answer is
cleanly extractable.

\begin{fromcode}[The specific generator]
The backbone is \code{Qwen-2.5-3B-Instruct}, a three-billion-parameter
instruction-tuned decoder, used through the chat-markup prompt format. Generation
is \emph{greedy} (\code{do\_sample=False}) for reproducibility, capped at $512$
new tokens for both the zero-shot and retrieval paths. Crucially, the backbone is
\textbf{frozen}: across the entire run its original weights never change. The only
thing that trains is a small adapter bolted onto it, which we cover in
\Cref{sec:prereq-lora}. Freezing the backbone is not an incidental choice; it is
the mechanism that bounds the catastrophic forgetting of \Cref{ch:problem}.
\end{fromcode}

\section{Embeddings: turning meaning into geometry}
\label{sec:prereq-embeddings}

The memory in \caem{} has to answer one question quickly: have we seen a query
like this one before? Comparing raw strings cannot do this, because ``Who wrote
\emph{Hamlet}?'' and ``Name the author of \emph{Hamlet}'' share almost no
characters yet mean the same thing. Embeddings solve this by mapping text to
points in a high-dimensional space, arranged so that similar meanings sit close
together.

\begin{intuition}[An embedding is a meaning-coordinate]
Imagine every possible sentence placed somewhere in a space with hundreds of
axes, positioned so that sentences with similar meaning land near one another and
unrelated sentences land far apart. An embedding model is the function that
computes those coordinates. Once two sentences are points, ``are they similar?''
becomes ``are these two points close?'', which is a geometry question a computer
answers in microseconds.
\end{intuition}

The standard closeness measure for embeddings is \emph{cosine similarity}, the
cosine of the angle between two vectors. It is $1$ when the vectors point the same
way, $0$ when they are orthogonal, and $-1$ when they are opposed. For two
vectors $\mathbf{a}$ and $\mathbf{b}$,
\[
  \cos(\mathbf{a}, \mathbf{b})
  \;=\;
  \frac{\mathbf{a} \cdot \mathbf{b}}{\lVert \mathbf{a}\rVert\,\lVert \mathbf{b}\rVert}.
\]
There is a useful trick here. If every embedding is \emph{normalised} to unit
length in advance, the denominator becomes one, and cosine similarity collapses
to a plain dot product. That single optimisation is why \caem{} normalises every
embedding the moment it is produced: it lets the similarity search use the fastest
possible inner-product machinery while still computing true cosine similarity.

\begin{figure}[htbp]
\centering
\begin{tikzpicture}[font=\small\sffamily, scale=1.0]
  % unit circle
  \draw[cbInk!25, line width=0.8pt] (0,0) circle (2.4);
  \draw[->, cbInk!50] (-2.9,0) -- (2.9,0);
  \draw[->, cbInk!50] (0,-2.9) -- (0,2.9);
  % vector a
  \draw[-{Stealth[length=2.4mm]}, line width=1.4pt, cbIntuition]
    (0,0) -- (2.4*0.866, 2.4*0.5)
    node[anchor=south west] {\,``Who wrote \emph{Hamlet}?''};
  % vector b (close)
  \draw[-{Stealth[length=2.4mm]}, line width=1.4pt, cbExample]
    (0,0) -- (2.4*0.985, 2.4*0.174)
    node[anchor=west] {\,``Name the author of \emph{Hamlet}''};
  % vector c (far)
  \draw[-{Stealth[length=2.4mm]}, line width=1.4pt, cbPitfall]
    (0,0) -- (2.4*0.0, 2.4*1.0)
    node[anchor=south] {``What is the boiling point of water?''};
  % angle arc between a and b (small)
  \draw[cbInk!55, line width=0.8pt] (1.2,0.69) arc (30:10:1.38);
  \node[font=\scriptsize, cbInk!70] at (1.95,0.55) {small angle};
  \node[font=\scriptsize, cbInk!70] at (0.55,1.55) {large angle};
  \node[font=\scriptsize, cbInk!60] at (0,-2.75) {(all vectors normalised to the unit circle)};
\end{tikzpicture}
\caption[Cosine similarity as geometry]{Cosine similarity is the cosine of the
angle between two normalised embeddings. The two paraphrases of the same question
(blue and green) point in nearly the same direction, so their cosine similarity is
high. The unrelated question (red) points elsewhere, so its similarity is low.
Once vectors are normalised to unit length, this cosine equals their dot product. (Schematic.)}
\label{fig:cosine}
\end{figure}

\begin{fromcode}[The specific embedding model]
The encoder is \code{sentence-transformers/all-mpnet-base-v2}, a sentence-level
embedding model that maps any text to a $768$-dimensional vector. \caem{} stores
the output as $32$-bit floats and L2-normalises every vector to unit length, so
that an inner product equals cosine similarity exactly. The code guards this
invariant aggressively: it validates the dimension on load (refusing a model that
does not produce $768$ dimensions) and re-normalises defensively if any vector's
length drifts from one. This same embedding is reused three times per query, for
the memory search, for grounding in the verifier, and for retrieval in the
retrieval-augmented tier.
\end{fromcode}

\section{Similarity search at scale: FAISS}
\label{sec:prereq-faiss}

Cosine similarity tells us how to compare \emph{two} vectors. But the memory may
hold up to a million episodes, and the retrieval corpus holds tens of millions of
passages. Comparing a query against every stored vector one by one would be far
too slow. This is the problem that a vector index solves, and \caem{} uses the
FAISS library for it.

There are two regimes, and \caem{} uses both. An \textbf{exact} index compares
the query against every vector but does so with heavily optimised inner-product
arithmetic; it returns the true nearest neighbour every time, and it is fast
enough while the collection is small. An \textbf{approximate} index trades a sliver
of accuracy for a large speed gain, and becomes necessary at scale. The
approximate scheme \caem{} uses combines two techniques:

\begin{description}[leftmargin=0pt, labelsep=0.5em, style=nextline]
  \item[Inverted file (IVF).] The vectors are clustered into a fixed number of
  cells, each with a representative centroid. At search time the query is compared
  only against the vectors in the few cells whose centroids are nearest, rather
  than the whole collection. The number of cells is the \emph{nlist} parameter and
  the number probed per query is \emph{nprobe}; together they trade recall against
  speed.
  \item[Product quantisation (PQ).] Each vector is split into sub-vectors, and
  each sub-vector is replaced by the index of the nearest entry in a small learned
  codebook. This compresses each vector to a handful of bytes, so that millions of
  vectors fit in memory and the distance arithmetic runs on the compact codes.
\end{description}

\begin{fromcode}[The two indexes inside \caem{}]
The episodic memory of \Cref{ch:memory} bootstraps as an \emph{exact}
inner-product index wrapped so that entries can be deleted by identifier, and then
\textbf{auto-promotes} to an approximate inverted-file product-quantised index once
it holds enough vectors to train the clustering (a threshold of twenty thousand
points; cells $=4096$, probes $=32$, sixty-four sub-quantisers at eight bits
each). Early cycles therefore enjoy exact search, and only large stores pay the
approximate-search tax. The retrieval corpus of \Cref{ch:tiers} is far larger and
is approximate from the start, with a finer clustering (cells $=32768$, probes
$=64$). The memory's novelty rule also rides on this search: a candidate episode
whose nearest stored neighbour exceeds a cosine similarity of $0.95$ is judged a
near-duplicate and is not stored again.
\end{fromcode}

\begin{bythenumbers}[The vector machinery at a glance]
\begin{itemize}[leftmargin=1.3em, itemsep=1pt]
  \item Embedding model: \code{all-mpnet-base-v2}, $768$ dimensions, unit-normalised.
  \item Episodic memory: exact index up to $20{,}000$ vectors, then inverted-file
  product-quantised; capacity $1{,}000{,}000$; novelty threshold $0.95$.
  \item Retrieval corpus: tens of millions of passages, approximate from the start.
  \item Similarity: cosine, computed as an inner product on normalised vectors.
\end{itemize}
\end{bythenumbers}

\section{Calibration: making a score mean something}
\label{sec:prereq-calibration}

This is the most important idea in the chapter, because it is the idea
\Cref{ch:problem} identified as the missing piece. A model can be accurate and yet
\emph{miscalibrated}: its confidence numbers do not match its actual hit rate. We
say a score is \emph{calibrated} when it can be read as a probability. If we
collect every answer a calibrated system rated $0.7$, then close to seventy
percent of them should be correct. Calibration is what lets a single fixed
threshold mean the same thing across very different questions.

\begin{intuition}[Calibrated versus merely high]
A confidence score is useful only if it is honest. A weather forecaster who says
``seventy percent chance of rain'' is well calibrated if it rains on about seventy
percent of such days, regardless of whether the forecaster is usually optimistic or
pessimistic. A language model, left alone, is a badly calibrated forecaster: it
says ninety percent and is right sixty percent of the time. Calibration is the
post-processing step that bends those raw scores back onto the line where the
number equals the frequency.
\end{intuition}

The standard way to measure the gap is the \emph{reliability diagram}: bucket
predictions by their stated confidence, and plot the average confidence in each
bucket against the actual accuracy in that bucket. A perfectly calibrated system
lies on the diagonal. The average vertical gap, weighted by bucket size, is the
\emph{expected calibration error}, and a curve sitting below the diagonal is the
signature of overconfidence.

\begin{figure}[htbp]
\centering
\begin{tikzpicture}
\begin{axis}[
  width=9.2cm, height=6.4cm,
  xlabel={stated confidence}, ylabel={actual accuracy},
  xmin=0, xmax=1, ymin=0, ymax=1,
  xtick={0,0.25,0.5,0.75,1}, ytick={0,0.25,0.5,0.75,1},
  axis lines=left, font=\small\sffamily,
  legend style={at={(0.03,0.97)}, anchor=north west, draw=cbInk!25,
    font=\scriptsize, fill=white},
  tick label style={font=\scriptsize},
  grid=both, grid style={cbInk!8},
]
  % perfect calibration diagonal
  \addplot[cbInk!55, dashed, line width=1pt, domain=0:1, samples=2] {x};
  \addlegendentry{perfectly calibrated}
  % overconfident curve (below diagonal)
  \addplot[cbPitfall, line width=1.4pt, mark=*, mark size=1.4pt] coordinates {
    (0.1,0.07) (0.25,0.16) (0.4,0.27) (0.55,0.39) (0.7,0.5) (0.85,0.62) (0.95,0.72)
  };
  \addlegendentry{raw model (overconfident)}
  % calibrated curve (close to diagonal)
  \addplot[cbExample, line width=1.4pt, mark=square*, mark size=1.3pt] coordinates {
    (0.1,0.11) (0.25,0.24) (0.4,0.41) (0.55,0.53) (0.7,0.69) (0.85,0.84) (0.95,0.94)
  };
  \addlegendentry{after calibration}
\end{axis}
\end{tikzpicture}
\caption[A reliability diagram]{A reliability diagram (schematic, drawn to illustrate
the concept). The raw model (red) sits well below the diagonal: when it says $0.85$
it is right only about $0.62$ of the time, the overconfidence of \Cref{ch:problem}.
Calibration (green) bends the scores back toward the diagonal, so a stated confidence
can be read as an actual probability. The vertical gap to the diagonal, averaged over
buckets, is the expected calibration error.}
\label{fig:reliability}
\end{figure}

\caem{} uses four calibration tools, each for a different signal, and it is worth
naming them once here so the later chapters can refer back.

\begin{description}[leftmargin=0pt, labelsep=0.5em, style=nextline]
  \item[Temperature scaling.] The simplest fix: divide the score's logit by a
  single learned temperature before the sigmoid. One number, fit to minimise the
  calibration error, slides the whole confidence curve toward the diagonal. \caem{}
  uses this for the pre-routing confidence of \Cref{ch:upre}, with a \emph{separate
  temperature per benchmark}.
  \item[Isotonic regression.] A more flexible, non-parametric fit that learns any
  monotone mapping from raw score to calibrated probability. It can bend the curve
  into any increasing shape, which matters when a signal is informative but not
  linearly so. \caem{} fits one isotonic curve \emph{per verifier signal} and
  combines them, as the next paragraph describes.
  \item[Platt scaling.] A logistic fit, two parameters, that puts the output of one
  model onto the probability scale of another. \caem{} uses it to align the
  long-context entailment judge with the short-context one (\Cref{sec:prereq-nli}),
  so that both judges report comparable numbers.
  \item[Shrinkage toward a pool.] When a per-benchmark fit is estimated from only a
  few hundred labelled examples, it overfits. The fix is to pull each per-benchmark
  curve partway toward the pooled curve estimated from all benchmarks at once.
  \caem{} sets this shrinkage weight to $0.6$, keeping moderate per-benchmark
  adaptation while regularising the weakest benchmarks.
\end{description}

\begin{precise}[How the nine signals become one calibrated number]
The verifier's composite is built so that the final score \emph{is} a calibrated
probability rather than an arbitrary weighted sum. Each of the underlying signals is
passed through its own isotonic curve, which converts the raw signal into a
calibrated per-signal probability and hence a log-odds contribution. The
contributions are summed, and a sigmoid turns the total back into a probability.
A signal that carries no information fits a nearly flat curve and contributes close
to zero log-odds automatically, with no manual zeroing. A signal whose useful
direction flips between benchmarks, for instance a relevance score that means
opposite things on a fact-checking task and an open-ended question, is handled
because its curve is fit per benchmark and can decrease on one and increase on
another. This is why a single pair of storage thresholds can govern five different
benchmarks: the comparison happens after every signal has been mapped onto the same
probability scale.
\end{precise}

\section{Natural-language inference: entailment as a grounding test}
\label{sec:prereq-nli}

The verifier's strongest evidence that an answer is grounded is not whether the
model felt confident, but whether the retrieved evidence actually \emph{supports}
the answer. The tool for that is natural-language inference.

Natural-language inference asks, given a \emph{premise} and a \emph{hypothesis},
whether the premise entails the hypothesis, contradicts it, or is neutral. \caem{}
puts the retrieved passage in the premise slot and the model's answer in the
hypothesis slot, and reads the entailment probability as a grounding score: a high
value means the evidence backs the answer, a low value means the answer floats free
of its evidence. This is the heart of the difference between a grounded answer and
a confabulation that merely sounds right.

\begin{pitfall}[Entailment is directional, and absence of support is not contradiction]
Two traps lurk here. First, entailment is \emph{directional}: ``the passage entails
the answer'' is the claim we want, not the reverse. Second, the two judges \caem{}
uses report only an entailment probability and treat \emph{not supported} the same
as \emph{neutral}; they do not raise a separate contradiction alarm. \caem{}
therefore relies on \emph{low grounding} as its low-support signal rather than a
hard contradiction veto. A reader who expects a three-way support, refute, neutral
verdict from these judges will misread the verifier; the operative signal is the
continuous grounding probability, not a discrete label.
\end{pitfall}

\begin{fromcode}[Two judges, dispatched by length]
\caem{} runs entailment through an adaptive dispatcher with two judges behind it.
A specialised, fast checker handles short hypotheses, and a long-context model
handles long ones. The split point is $408$ tokens, and it is not arbitrary: the
fast checker's encoder caps at $512$ tokens, of which $4$ are reserved for prompt
overhead and $100$ for a minimum slice of premise, leaving $512 - 4 - 100 = 408$
tokens of hypothesis that fit without truncation. Anything longer routes to the
long-context judge, whose output is Platt-scaled (\Cref{sec:prereq-calibration})
against the fast checker on a shared five-hundred-sample fold, so the two report
numbers on the same scale. The dispatcher is deliberately invisible to the rest of
the verifier, which sees one entailment probability and never needs to know which
judge produced it. One honest caveat applies to the deployed run: the long-context
judge was set aside because its post-hoc probability fit did not clear the required
correlation floor, so the shipped verifier ran the fast checker for every
prediction, with longer hypotheses truncated at the checker's window (a limitation
recorded among the threats).
\end{fromcode}

\section{Low-rank adaptation: learning without forgetting}
\label{sec:prereq-lora}

The last tool is the one that powers the self-improvement loop. The challenge from
\Cref{ch:problem} was that retraining a model on new material tends to erase its
old capabilities. Low-rank adaptation, or LoRA, is the technique that lets \caem{}
retrain on its own verified answers each cycle while bounding that risk.

\begin{intuition}[Freeze the library, add a thin notebook]
Picture the model's knowledge as a vast reference library that took enormous effort
to assemble. Naive fine-tuning rewrites pages of that library to absorb new facts,
and in doing so smudges unrelated pages, which is catastrophic forgetting. LoRA
instead leaves the library completely untouched and attaches a thin notebook beside
it. The notebook is small, it is the only thing that gets written, and at any time
it can be torn out and the library is exactly as it was. \caem{} writes a fresh
page into that notebook each cycle, and the retention guard of \Cref{ch:cycle} is
the rule that tears the page out if it made things worse.
\end{intuition}

The mechanism is a low-rank matrix. A weight matrix $W$ inside the frozen model maps
an input $\mathbf{x}$ to an output $W\mathbf{x}$. LoRA freezes $W$ and adds a
trainable correction built from two small matrices, $A$ and $B$, whose product has a
deliberately low rank $r$:
\[
  \mathbf{h} \;=\; W\mathbf{x} \;+\; \frac{\alpha}{r}\,B A\,\mathbf{x},
  \qquad
  A \in \mathbb{R}^{r \times k},\;
  B \in \mathbb{R}^{d \times r},\;
  r \ll \min(d, k).
\]
Only $A$ and $B$ are learned; $W$ never moves. The scalar $\alpha/r$ controls how
strongly the correction is applied. Because $r$ is tiny, the number of trainable
parameters is a small fraction of the matrix it corrects.

\begin{workedexample}[Why ``low rank'' is such a large saving]
Take a single square projection of size $d \times d$. Training it fully means
learning $d^2$ numbers. LoRA instead learns $A$ and $B$ with $r\,d$ numbers each,
for $2rd$ in total. With a rank of $r=32$ and a typical projection a couple of
thousand wide, $2rd$ is on the order of a few percent of $d^2$, roughly a thirty-fold
reduction \emph{on that one matrix}. Aggregated over the handful of projections
\caem{} adapts, the trainable footprint comes to roughly two percent of the full
model. The other ninety-eight percent stays frozen, which is exactly the property
that keeps a cycle's update from rewriting unrelated knowledge.
\end{workedexample}

\begin{fromcode}[The specific adapter]
\caem{} adapts with rank $r=32$ and scaling $\alpha=64$ (the standard ``alpha equals
twice the rank'' rule, giving a correction scale of two), a dropout of $0.05$, and a
learning rate of $2\times 10^{-4}$. The adapter is attached to all seven linear
projections of each transformer block: the four attention projections (query, key,
value, output) and the three feed-forward projections (gate, up, down). This is the
``adapt every linear layer'' recommendation from the LoRA-scaling literature.
Because the base is frozen, the explicit anchor term that a full fine-tune would
need to stay near its starting point is dropped on the LoRA path; the freezing does
that job for free. This adapter is the only thing the self-improvement loop of
\Cref{ch:cycle} ever changes, and rolling a cycle back means nothing more than
discarding the latest version of these two small matrices per projection.
\end{fromcode}

\section{How the toolbox maps onto the architecture}
\label{sec:prereq-map}

Each tool reappears at a definite place, and it helps to fix the correspondence
before moving on. \Cref{tab:prereq-map} is that correspondence.

\begin{table}[htbp]
\centering
\caption[Concept-to-stage map]{Where each prerequisite concept does its work in the
chapters ahead.}
\label{tab:prereq-map}
\small
\renewcommand{\arraystretch}{1.25}
\begin{tabularx}{\textwidth}{@{}l >{\raggedright\arraybackslash}X l@{}}
\toprule
\textbf{Concept} & \textbf{Where it does its work in \caem{}} & \textbf{Chapter} \\
\midrule
Per-token probability & The $\utoken$ signal in pre-routing and verification
  & \Cref{ch:upre,ch:verifier} \\
Embeddings + cosine & Encoding queries and matching them to stored episodes
  & \Cref{ch:memory,ch:router} \\
FAISS search & Memory lookup and retrieval-corpus search at scale
  & \Cref{ch:memory,ch:tiers} \\
Calibration & Honest pre-routing confidence and the calibrated composite
  & \Cref{ch:upre,ch:composite} \\
Entailment (NLI) & Grounding the answer against retrieved evidence
  & \Cref{ch:verifier} \\
Low-rank adaptation & Learning from verified answers without forgetting
  & \Cref{ch:cycle} \\
\bottomrule
\end{tabularx}
\end{table}

\begin{defence}[If an examiner asks: ``which of these tools are novel?'']
None of the six tools is novel, and the thesis does not claim them. Embeddings,
similarity search, calibration, entailment, and low-rank adaptation are all standard
machinery. The contribution of \caem{} is not a new tool but the \emph{architecture}
that composes them: a calibrated confidence deciding how hard to work, an
entailment-grounded composite gating a four-outcome admission rule, and a low-rank
update under a retention guard closing the loop. The novelty is in the wiring, and
the wiring is what the rest of the book is about.
\end{defence}

With the toolbox assembled, we are ready to open the first box. \Cref{ch:upre}
begins the per-query pipeline with the pre-routing confidence, the single number
that decides, before any answer is generated, how much work a query deserves.
